% META: {"id": "TNSI-STRUCT-CO-049", "chapitre": "TNSI-STRUCTURES-DONNEES", "type_objet": "correction", "exercice_ref": "TNSI-STRUCT-EX-001", "capacites_codes": ["C14"], "sources_inspiration": [], "status": "approved", "capacites": ["TNSI-STRUCTURES-DONNEES-C14"], "parcours": 2, "fichier_exercice": "chapitres/TNSI-STRUCTURES-DONNEES/exercices/TNSI-STRUCT-EX-049.tex"}
% BEGIN-VERIFY
% class Pile:
%     def __init__(self):
%         self._d = []
%     def empiler(self, v):
%         self._d.append(v)
%     def depiler(self):
%         return self._d.pop()
%     def est_vide(self):
%         return len(self._d) == 0
% def est_equilibre(expr):
%     p = Pile()
%     paires = {')': '(', ']': '[', '}': '{'}
%     for c in expr:
%         if c in '([{':
%             p.empiler(c)
%         elif c in ')]}':
%             if p.est_vide() or p.depiler() != paires[c]:
%                 return False
%     return p.est_vide()
% assert est_equilibre('([{}])') == True
% assert est_equilibre('([)]') == False
% END-VERIFY
\begin{corrige}{TNSI-STRUCT-EX-001}

\textbf{Q1.} Une pile est adaptee car on doit toujours refermer le symbole ouvert le \textbf{plus recemment} en premier (comportement LIFO) : chaque symbole ouvrant est empile, et chaque symbole fermant doit correspondre au sommet de la pile.

\textbf{Q2.}
\begin{python}
class Pile:
    def __init__(self):
        self._d = []

    def empiler(self, v):
        self._d.append(v)

    def depiler(self):
        return self._d.pop()

    def est_vide(self):
        return len(self._d) == 0


def est_equilibre(expr):
    p = Pile()
    paires = {')': '(', ']': '[', '}': '{'}
    for c in expr:
        if c in '([{':
            p.empiler(c)
        elif c in ')]}':
            if p.est_vide() or p.depiler() != paires[c]:
                return False
    return p.est_vide()
\end{python}

\textbf{Q3.} \lstinline|est_equilibre("([{}])")| renvoie \lstinline{True} ; \lstinline{est_equilibre("([)]")} renvoie \lstinline{False} (le \lstinline{)} rencontre alors que le sommet de la pile est \lstinline{[}, qui ne correspond pas).

\end{corrige}
